\section{Related work}
\label{sec:related}

\paragraph{Sparse activations and training.}
Replacing smooth MLP nonlinearities with ReLU can yield high activation
sparsity while preserving language-model quality \citep{mirzadeh2023relu}.
ProSparse stages regularization and thresholding during pretraining
\citep{song2024prosparse}; the Sparsing Law studies scaling behavior under
sparsity-inducing training \citep{liu2024sparsing}; and recent work examines
sparsity-aware pretraining recipes at larger scales \citep{cetin2026sparser}.
Conditional-computation and activation-sparsity studies likewise motivate
training models whose dense weights are selectively exercised
\citep{zhang2024exploring}.  Our study is smaller and diagnostic: it holds a
one-pass data budget fixed, reports exact product counts, and does not claim a
compute-optimal sparse training recipe.

\paragraph{Post-training activation sparsification.}
TEAL selects activation thresholds without retraining and evaluates a
model-level, matrix-footprint-weighted configuration \citep{liu2024teal}.
CATS and CHESS introduce contextual or structured activation selection
\citep{lee2024cats,zhang2024chess}, while R-Sparse combines rotations and
thresholds to increase inference-time activation sparsity
\citep{sun2025rsparse}.  These works make efficiency concrete through kernels
or latency measurements.  We use a TEAL-like evaluation-only clipping family
as a reference, but our central object is different: zero-containing scalar
products reconstructed from actual operands, including causal attention reuse
and a dense output head.  Our metric is useful before kernel engineering, but
cannot replace it.

\paragraph{Gradient interaction.}
Projecting conflicting auxiliary gradients has precedents in multi-task
optimization, including PCGrad \citep{yu2020pcgrad} and Adam-aware approaches
such as Bloop \citep{mitrevski2024bloop}.  Our orthogonal $\ell_1$ update uses
the task optimizer's second moment and a relative direction budget.  It enforces
a local non-opposition condition on the reconstructed directions; it does not
guarantee that a finite parameter step improves task loss.  The present
comparison also does not isolate this optimizer from every topology and gate
choice, so we treat it as part of a complete recipe rather than a standalone
algorithmic contribution.

